A \code{BPatch\_frame} object represents a stack frame. The
getCallStack member function of \code{BPatch\_thread} returns a vector of
\code{BPatch\_frame} objects representing the frames currently on the stack.

\begin{apient}
  enum BPatch_frameType {
    BPatch_frameNormal,     // Normal stack frame
    BPatch_frameSignal,     // Signal invocation
    BPatch_frameTrampoline, // Call into instrumentation code
  }
\end{apient}

\begin{apient}
  BPatch_frameType getFrameType()
\end{apient}

\apidesc{
  Return the type of the stack frame.
}

\begin{apient}
  void *getFP()
\end{apient}

\apidesc{
  Return the frame pointer for the stack frame.
}
\begin{apient}
  void *getPC()
\end{apient}

\apidesc{
  Returns the program counter associated with the stack frame.
}
\begin{apient}
  BPatch_function *findFunction()
\end{apient}

\apidesc{
  Returns the function associated with the stack frame.
}
\begin{apient}
  BPatch_thread *getThread()
\end{apient}

\apidesc{
  Returns the thread associated with the stack frame.
}
\begin{apient}
  BPatch_point *getPoint()
  BPatch_point *findPoint()
\end{apient}

\apidesc{
  For stack frames corresponding to inserted instrumentation, returns
  the instrumentation point where that instrumentation.
}

\begin{apient}
  bool isSynthesized()
\end{apient}

\apidesc{
  Returns true if this frame was artificially created, false otherwise.
}
